\subsection{Module D.7: disk--finite-set calipers}

\paragraph{Exact interface.}
The enclosure objective is $\Lambda(\mathcal D_\eta\cup P)$ with $\Lambda$
and the closed centered disk defined in
\zcref{def:finite-enclosure-number}; here $P$ is the displayed finite set.

\paragraph{Variables and feasible set.}
Let $\eta\ge0$, let $\mathcal D_\eta=\mathbb B(O,\eta)$, and take
$P=\{p_1,\ldots,p_m\}$.  The feasible orientation is a unit normal on the
full normal circle; point--disk ties require $\lVert p_i\rVert\ge\eta$, and
point--point ties require $p_i\ne p_k$.

\paragraph{Objective.}
Minimize the three-direction support sum for $\mathcal D_\eta\cup P$ and
reduce every non-disk-only minimizer to a finite list of tie normals.

\paragraph{Exact result.}

\begin{proposition}[Disk--finite-set calipers]
\label{prop:app-disk-finite-caliper}
Let $\mathcal D_\eta$ be centered at $O$, let
$P=\{p_1,\ldots,p_m\}$, and let $\mathsf J$ denote rotation by $\pi/2$.
Unless the disk is the active support in all three directions at a minimizing
orientation, a minimizing normal belongs to the finite list
\[
 \pm\frac{\mathsf J(p_k-p_i)}{\|p_k-p_i\|}
 \quad(p_i\ne p_k)
\]
or
\[
 \frac{\eta p_i\pm
 \sqrt{\|p_i\|^2-\eta^2}\,\mathsf Jp_i}{\|p_i\|^2}
 \quad(\|p_i\|\ge\eta).
\]
Thus one side of a least enclosing equilateral triangle either contains two
points, or is tangent to the disk and contains one point.
\end{proposition}

\paragraph{Solution.}

\begin{proof}[Proof of \zcref{prop:app-disk-finite-caliper}]
The tie directions partition the normal circle into cells with fixed support
sources.  On a cell with $q$ disk supports,
\[
 \sum_{j=0}^2h_{\mathcal D_\eta\cup P}(\mathsf R^jn(\theta))
 =
 q\eta+\langle v,n(\theta)\rangle.
\]
If any point is active, the last inner product is positive and the second
derivative is its negative.  Hence an interior critical point is a maximum.
A minimum occurs at a point--point or point--disk tie.  On a disk-only cell
the support sum is $3\eta$, already the global minimum.
\end{proof}

\paragraph{Body consequence.}
This finite caliper reduction supplies the disk-plus-point evaluation used
by \zcref{thm:new-complementary-gap}.

\subsection{Module D.8: half-edge one-third radial envelope}

\paragraph{Exact interface.}
Use the anchor admissible set and attained radial envelope $c_{\max}$ from
\zcref{def:finite-anchor-admissible-set}.

\paragraph{Variables and feasible set.}
The ordered boundary demands satisfy
$M\ge1/2$, $0<m\le M$, and $M+m<1$.  The candidate radial coordinate is in
$[0,1]$ and is selected from the $L$ or $T$ cell of
\zcref{lem:app-local-admissible-set}.

\paragraph{Objective.}
Give a strict upper bound for the maximum feasible own-radial reach in terms
of the smaller boundary demand.

\paragraph{Exact result.}

\begin{lemma}[Half-edge one-third radial envelope]
\label{lem:app-one-third-radial-envelope}
If
\[
 M\ge\frac12,\qquad 0<m\le M,\qquad M+m<1,
\]
then
\[
 \boxed{c_{\max}(M,m)<1-\frac m3.}
\]
\end{lemma}

\paragraph{Solution.}

\begin{proof}[Proof of \zcref{lem:app-one-third-radial-envelope}]
Put $s=M+m$ and $c_0=1-m/3$.  Only the $L$ and $T$ cells of
\zcref{lem:app-local-admissible-set} can occur.

For $m\le3/8$, substitution into the $L$ polynomial gives
\[
 F_L(c_0)
 =
 \frac m{81}(m^3-12m^2-63m+27)>0;
\]
the cubic decreases on $[0,3/8]$ and equals $891/512$ at $3/8$.
Moreover
\[
 F_L\!\left(\frac{\sqrt3}{2}\right)
 =-\left(m-\frac{\sqrt3}{4}\right)^2\le0,
\]
and
\[
 F_L'(c)=c(4c^2-2)+m>0
 \qquad\left(c\ge\frac{\sqrt3}{2}\right).
\]
Thus $F_L$ is strictly increasing from its selected sign change on
$[\sqrt3/2,1]$, and the selected $L$ root is below $c_0$.  At the $L/T$
transition the two roots agree.  On the $T$ cell,
\[
 c_T(s)=\frac{2(s-m)}{1+\sqrt{4s^2-3}},
\]
and, writing $r=\sqrt{4s^2-3}$,
\[
 c_T'(s)=\frac{2(r+4ms-3)}{r(1+r)^2}<0
\]
because $m\le s/2$ and $r+2s^2-3<0$ for $s<1$.  Hence the $T$ root remains
below $c_0$.

For $3/8\le m<1/2$, one has
\[
 s^4-s^2+mM
 \ge
 \left(\frac78\right)^4-\left(\frac78\right)^2+\frac3{16}
 =\frac{33}{4096}>0,
\]
so only the $T$ cell occurs.  It decreases with $M$; at $M=1/2$,
$s\ge7/8$ and its value is at most $4/5<5/6<c_0$.
\end{proof}

\paragraph{Body consequence.}
The strict one-third bound supplies the local endpoint estimate in the
adjacent branch of \zcref{prop:finite-vd-radial-separation}.

\subsection{Module D.9: rescuer endpoint and boundary-tail budget}

\paragraph{Exact interface.}
Use the supported interval, actual reaches, strict-supercritical envelope,
and four-edge nonsupercritical path in
\zcref{prop:new-one-t3-terminal,prop:finite-vd1-supported-terminal}, together
with the exact signed center traces and exits of
\zcref{prop:signed-center-normal-form}.

\paragraph{Variables and feasible set.}
The special role has boundary endpoint $a$ and radial interval
$[c,u]\subset r_1$ with $\varepsilon=1-u$ and
$M=M_c^{\mathrm{sup}}$.  The center role is CE1 or CE2 in signed normal form,
with $0<R<1$, $W=1-R$, $E=\sqrt{1-RW}$,
$\eta=1-E$, positive slacks $\alpha,\delta$, and
$k=\eta+\alpha+\delta$.  The four displayed adapter inequalities are
assumed, $T_1$ is uniquely supercritical with radial reach at least $c$, and
$T_2,\ldots,T_5$ form the center-free nonsupercritical path.

\paragraph{Objective.}
Lower-bound the far boundary demand by $M$ in both center-trace alternatives
and test it against the four nonsupercritical boundary budgets.

\paragraph{Exact result.}

\begin{lemma}[Rescuer endpoint and boundary-tail budget]
\label{lem:app-rescuer-tail-budget}
Let $T_C$ be a CE1 or CE2 center role in the signed normal form of
\zcref{prop:signed-center-normal-form}.  Thus
\[
 0<R<1,
 \quad W=1-R,
 \quad E=\sqrt{1-RW},
 \quad \eta=1-E,
 \quad \alpha,\delta>0,
 \quad k=\eta+\alpha+\delta,
\]
when its positive companion boundary trace on $e_{5,0}$ is present, that
trace is $[k/W,R+\alpha]$; it may instead be empty or a point contact.  Its
exit on $r_1$, measured from $O$, is $\delta/R$.
Suppose a special role at $V_0$ has boundary endpoint $a$ on $e_{5,0}$ and
supported interval $[c,u]\subset r_1$, measured from $V_1$ toward $O$.
Put
\[
 \varepsilon=1-u,\qquad M=M_c^{\rm sup}.
\]
Assume
\[
 \varepsilon V_1\in U_C,\qquad a+\varepsilon\le1,
\]
\[
 a\le1-M,\qquad
 \frac{a}{a+\varepsilon}\le1-M.
\]
Assume also
\[
 U_C\cap\bigcup_{i=1}^4e_{i,i+1}=\varnothing.
\]
If $T_1$ is uniquely supercritical with own-radial reach at least $c$ and
$T_2,T_3,T_4,T_5$ are nonsupercritical along this four-edge path,
then the skeleton is not covered.
\end{lemma}

\paragraph{Solution.}

\begin{proof}[Proof of \zcref{lem:app-rescuer-tail-budget}]
Let $h_T$ be the far boundary demand on $T_5$.  If the companion C trace
does not hide $a$, then $h_T\ge1-a\ge M$.

In the hiding configuration, $k/W\le a$, while
$\varepsilon V_1\in U_C$ gives $\delta/R\ge\varepsilon$.  Therefore
\[
 Wa\ge k=\eta+\alpha+\delta>\alpha+R\varepsilon,
\]
so
\[
 a-R(a+\varepsilon)>\alpha.
\]
Also $Wa>\delta\ge R\varepsilon$, so
$R<a/(a+\varepsilon)$.  Since $a+\varepsilon\le1$,
\[
 R+\alpha
 <a+R(1-a-\varepsilon)
 <\frac{a}{a+\varepsilon}
 \le1-M.
\]
Thus again $h_T\ge M$.

The strict-supercritical envelope gives $B_1<M$.  Adding
\[
 A_2\ge1-B_1,\quad
 B_2+A_3\ge1,\quad B_3+A_4\ge1,\quad B_4+A_5\ge1,
 \quad B_5\ge h_T
\]
gives
\[
 \sum_{i=2}^5(A_i+B_i)>4,
\]
contrary to nonsupercriticality.
\end{proof}

\paragraph{Body consequence.}
Together with the two branch-specific adapters, this common terminal proves
\zcref{prop:new-one-t3-terminal,prop:finite-vd1-supported-terminal}.
